\documentclass[12pt]{article}
\usepackage{amsmath}
\usepackage{amssymb}

\begin{document}

\section*{Question 1}

Consider a long shell of radius $R$ with a uniform surface charge density $\sigma$. Inside this cylindrical shell, coaxially placed, is a long solid cylinder of radius $a$ ($a<R$) with a uniform volume charge density $\rho$.

\begin{enumerate}
\item[(a)] Find the electric field in the regions $0 \leqslant r \leq a$, $a \leqslant r \leq R$ and $r \geqslant R$ -- 0.6 pts

\item[(b)] Determine the potential difference $V$ between the surface of the solid cylinder (radius $a$) and the cylindrical shell (radius $R$) -- 0.4 pts
\end{enumerate}

\section*{Question 2}

Consider a long, straight cylindrical conductor with radius $R$. The conductor carries a current $I$ uniformly distributed across its cross-sectional area. Surrounding the conductor is a cylindrical shell with an inner radius $a$ (where $a>R$) and outer radius $b$ (where $b>a$). The shell carries a current $-I$ uniformly distributed across its cross-sectional area in the opposite direction.

\begin{enumerate}
\item[(a)] Find the magnetic field at a distance $r$ from the axis of the conductor for the following regions: $0 \leqslant r \leq R$, $R \leqslant r \leq a$, $a \leqslant r \leq b$ and $r \geqslant b$ -- 0.8 pts

\item[(b)] Sketch the magnetic field as a function of $r$ -- 0.2 pts
\end{enumerate}

\section*{Question 3}

Consider a circular loop of radius $R$ made of a conducting wire, lying in the $xy$-plane with its center at the origin. The loop is placed in a time-varying magnetic field given by $\mathbf{B}(t)=B_0\left(\frac{t}{\tau}\right)\hat{\mathbf{z}}$, where $B_0$ is a constant magnetic field strength, $\tau$ is a characteristic time, and $t$ is time.

\begin{enumerate}
\item[(a)] Calculate the magnetic flux through the loop as a function of time -- 0.2 pts

\item[(b)] Determine the induced electromotive force (emf) in the loop as a function of time -- 0.2 pts

\item[(c)] Determine the direction of the induced current in the loop -- 0.2 pts

\item[(d)] Calculate the induced current in the loop as a function of time given that the resistance of the loop is $R$ -- 0.2 pts

\item[(e)] If there is an additional constant magnetic field directed along the $y$-axis, $\mathbf{B}_c=B_y\hat{\mathbf{y}}$, what would be the magnitude and direction of the torque acting on the loop? -- 0.2 pts
\end{enumerate}

\section*{Question 4}

Consider an electromagnetic wave propagating in vacuum with electric field $\mathbf{E}$ and magnetic field $\mathbf{B}$. The wave is described by the following equations:
\begin{equation*}
\begin{split}
    \mathbf{E}(z,t)&=E_0\cos(kz-\omega t)\hat{x} \\
    \mathbf{B}(z,t)&=B_0\cos(kz-\omega t)\hat{y}
\end{split}
\end{equation*}
where $E_0$ and $B_0$ are the amplitudes of the electric and magnetic fields respectively, $k$ is the wave number, and $\omega$ is the angular frequency. Here $c=\omega/k$, where $c$ is the speed of light.

Show that these solutions satisfy Maxwell's equations in vacuum and/or show how $E_0$ and $B_0$ must be related in order to satisfy them.

\end{document}
